%%% SECTION 4: ADAPTED ICH E6(R3) FOR PHYSICAL AI GOOD CLINICAL PRACTICE %%%

[This section should be approximately 10-12 pages. The major theme is that the
Adaption: ICH Harmonised Guideline \cite{ich-adapt} establishes comprehensive
good clinical practice principles specifically designed for Physical AI oncology
trials, building on the internationally recognized ICH E6(R3) framework
\cite{ich-e6r3}. By adapting an existing standard rather than creating a new
one, this work adds credibility to Physical AI implementations through
association with decades of established GCP practice.]

[PRIMARY SOURCE: national-platform/ich\_e6r3\_adapt/ directory containing four
chunked files: 01\_preamble\_principles\_investigator.tex (lines 1-337),
02\_sponsor\_responsibilities.tex (lines 338-683),
03\_data\_governance.tex (lines 684-911), and
04\_appendices\_glossary.tex (lines 912-1300). All four files must be used
together to build this section.]

\subsection{Principles of Physical AI Clinical Practice}
\label{subsec:ich-principles}

[Build from national-platform/ich\_e6r3\_adapt/01\_preamble\_principles\_investigator.tex,
Section 1 (subsections 1.1-1.5). Cover the adapted principles that govern
Physical AI clinical practice, including how traditional GCP principles
(beneficence, justice, respect for persons) are extended to include robotic
safety, AI transparency, and Physical AI system validation requirements.]

[MAJOR THEME: These principles establish that Physical AI systems must meet the
same ethical and quality standards as human-directed clinical practice, with
additional requirements for system validation, safety monitoring, and
autonomous decision documentation. This demonstrates that Physical AI does not
lower the bar for clinical practice but rather raises it with additional
safeguards.]

\subsection{Investigator Responsibilities in Physical AI Trials}
\label{subsec:ich-investigator}

[Build from national-platform/ich\_e6r3\_adapt/01\_preamble\_principles\_investigator.tex,
Section 2 (subsections 2.1-2.12). Cover the twelve categories of investigator
responsibilities as adapted for Physical AI contexts, including:
qualifications and training for Physical AI system oversight, adequate
resources including robot maintenance, medical care with Physical AI
integration, communication with IRB/IEC regarding Physical AI systems,
compliance with protocol including robot operational parameters, and
informed consent procedures specific to Physical AI interactions.]

[Explain how investigator responsibilities shift when Physical AI systems
perform functions traditionally done by human staff. MAJOR THEME: More patients
will be treated with more robots and fewer workers, but this shift requires
investigators to develop new competencies in Physical AI oversight.]

\subsection{Sponsor Responsibilities in Physical AI Trials}
\label{subsec:ich-sponsor}

[Build from national-platform/ich\_e6r3\_adapt/02\_sponsor\_responsibilities.tex,
Section 3 (subsections 3.1-3.15). Cover quality management systems for
Physical AI, regulatory submission requirements, IRB/IEC review of Physical AI
systems, trial design incorporating robotic workflows, monitoring Physical AI
system performance, safety reporting for Physical AI adverse events, data
handling with robot-generated records, and clinical trial reports documenting
Physical AI outcomes.]

[MAJOR THEME: Sponsors bear primary responsibility for ensuring Physical AI
systems are safe, validated, and properly integrated into trial operations.
The adapted ICH standard provides a comprehensive framework that builds on
existing sponsor obligations while adding Physical AI-specific requirements.
Reference the 24-hour simulation from Clinical Trial Site Complete
Documentation \cite{site-docs} as evidence that these responsibilities can be
fulfilled at scale.]

\subsection{Data Governance for Physical AI Trials}
\label{subsec:ich-data-governance}

[Build from national-platform/ich\_e6r3\_adapt/03\_data\_governance.tex,
Section 4 (subsections 4.1-4.3). Cover blinding safeguards in Physical AI
contexts (how to maintain blinding when robots are visible to participants),
data lifecycle elements for robot-generated data (telemetry, sensor readings,
autonomous decision logs, imaging data), and computerized systems requirements
for Physical AI platforms.]

[Reference FHIR \cite{hl7-fhir}, DICOM \cite{dicom}, and MCP
\cite{mcp-protocol} as interoperability standards that support Physical AI
data governance. Connect to the Physical AI National MCP Servers
\cite{national-mcp-paper} as the infrastructure implementation of these data
governance requirements.]

\subsection{Physical AI System Documentation}
\label{subsec:ich-documentation}

[Build from national-platform/ich\_e6r3\_adapt/04\_appendices\_glossary.tex,
Appendix A (Physical AI System Documentation). Cover the documentation
requirements for Physical AI systems in clinical trials, including system
specifications, validation reports, safety assessments, maintenance logs, and
software version control. Explain how these documentation requirements
ensure traceability and accountability for all Physical AI operations.]

\subsection{Clinical Trial Protocol Adaptations}
\label{subsec:ich-protocol}

[Build from national-platform/ich\_e6r3\_adapt/04\_appendices\_glossary.tex,
Appendix B (Clinical Trial Protocol). Cover how clinical trial protocols
must be adapted to include Physical AI system specifications, robot
operational parameters, fallback procedures for robot failures, and
criteria for when human staff must override Physical AI decisions.]

\subsection{Essential Records for Physical AI Trials}
\label{subsec:ich-records}

[Build from national-platform/ich\_e6r3\_adapt/04\_appendices\_glossary.tex,
Appendix C (Essential Records). Cover the additional essential records
required for Physical AI trials beyond standard GCP documentation,
including robot calibration records, AI model version histories, sensor
validation logs, and Physical AI adverse event reports.]

\subsection{Physical AI Glossary and Definitions}
\label{subsec:ich-glossary}

[Build from national-platform/ich\_e6r3\_adapt/04\_appendices\_glossary.tex,
Glossary section. Include key Physical AI definitions that are specific to
the ICH adaptation. Cross-reference with the 30 Physical AI-specific
definitions mentioned in the ICH adaptation. Note how these definitions
provide standardized terminology for the entire National Platform.]
